% =============================================================================
\chapter{16 September: what the cells mean, T3, and a bias in our own table}\label{ch:tables}

\begin{ideia}[The chapter in two sentences]
A cell of the price table answers one question: forced to cover honestly, how wide would this
procedure have to be, compared with discarding the knowledge? And the table of named reports, read
in full, shows that one retraction usually returns the analyst to the whole class, while the rule
that chose its rows biased the one column that looked most like Rosenbaum.
\end{ideia}

Two documents were written this morning, both in English and both on the branch:
\nolinkurl{rho_breakdown/table1_dossier_2026-09-16/}, a dossier of every Table~1 the results support,
and \nolinkurl{rho_breakdown/panelB_T3_2026-09-16/}, which opens one cell of the price table, prints the
named reports in full and audits how their rows were chosen.

\begin{armadilha}[Three numbering schemes]
The dossier's \textbf{T3} is M1's Table~2 (\texttt{results/interval/TABLE2\_REPORTS.md}), the paper's
future Table~2. The ``Panel~B'' of its recommended \textbf{T11} is the price table on M1's
\emph{Panel~A} (semi-synthetic), not M1's Panel~B. ``Table~3'' in the outline of 13 September is the
audit under wrong knowledge. Below, T-numbers are the dossier's and panels are M1's.
\end{armadilha}

\begin{antes}
\antesq[metrica-unica]{In the one-wrong-claim cell, what are $\lambda^{*}$ and the mean half-width of
the committed report, and how do they give 5.123?}
\antesq[metrica-unica]{How does a procedure earn an $\Lnf$ above 1?}
\end{antes}

\section{Opening one cell}

Take the one-wrong-claim cell of Panel~A at $n = 1000$: 359 queries, twenty seeds each, 7,180
replicates. For each replicate $i$ a procedure reports an interval with midpoint $m_i$ and half-width
$h_i$, and $\lambda_i = |\tau - m_i| / h_i$ is the factor by which that interval must be stretched about
its own midpoint to reach the true effect. Below 1, it already covers. Take $\lambda^{*}$, the 95th
percentile of the $\lambda_i$ over the cell: stretching every interval by $\lambda^{*}$ makes 95\,\% of
them cover. Average the stretched widths $2\lambda^{*}h_i$ and divide by the same quantity for the
whole class.

\begin{medido}[The cell, recomputed from replicates by \texttt{make\_figs.py}]
Script \nolinkurl{rho_breakdown/panelB_T3_2026-09-16/make_figs.py}, reading
\texttt{results/interval/replicates.csv.gz}. Panel~A, $b = 1$, $n = 1000$, 7,180 replicates.

\medskip
\begin{tabular}{@{}lrrrr@{}}
procedure & $\lambda^{*}$ & mean half-width & calibrated width & $\Lnf$ \\
\midrule
PLAIN, the committed report & 20.416 & 0.3855 & 15.7397 & \textbf{5.123} \\
$\Ir{1}$, one revision       & 1.000  & 1.3064 & 2.6129  & \textbf{0.851} \\
BLANKET, the whole class     & 1.000  & 1.5361 & 3.0721  & 1.000 \\
\end{tabular}

\medskip
$15.7397 / 3.0721 = 5.123$ and $2.6129 / 3.0721 = 0.851$. At $n = 20\,000$ the same two cells read
14.368 and 0.758: more data narrows a wrongly centred interval without moving it.
\end{medido}

The committed report is the narrowest interval in the cell (mean half-width 0.39 against 1.54 for the
class) and the most expensive once it is required to be right. \figref{fig:t3-why} shows why: the
stretch is about the interval's own midpoint, so a narrow interval centred on the effect the wrong
graph implies has to be compensated on both sides.

\begin{figure}[!ht]
\centering
\includegraphics[width=0.74\linewidth]{t3_f5_why_above_one.pdf}
\caption{Schematic. A narrow interval centred away from the truth needs a larger factor than a wide
honest one; that is the only way a procedure lands above 1, wider than discarding the knowledge.}
\label{fig:t3-why}
\end{figure}

\textbf{The reproduction.} The script reimplements the $\Lnf$ pipeline from \texttt{replicates.csv.gz}
and \texttt{rows.csv} and asserts that all 27 Panel~A cells of \texttt{TABLE1\_L95.md} at $n = 1000$
(nine procedures, three columns) come out to three decimals; I reran it today and it passes. The filter
that makes it reproduce: a query enters a column iff it is informative \emph{and} its realised $b$ is
non-empty (391, 359, 131 queries).

\begin{armadilha}[\texttt{profiles.jsonl} is the wrong join key for cells]
My first attempt joined the replicates on \texttt{results/interval/profiles.jsonl} and got 0.180 for
PLAIN at $b = 0$ instead of 0.464. That file holds 934 profiles, only for queries whose analyst graph
commits to a set, so the join silently rebuilt the committed-set secondary table, whose cell really is
0.180 on 134 queries. A correct number for the wrong population; the comparison with the published
cell is what caught it.
\end{armadilha}

\prompt[conceito=metrica-unica,dificuldade=0.7]{In the Panel~A, $b = 1$, $n = 1000$ cell, give
$\lambda^{*}$ and the mean half-width of the committed report and of the class, and compute $\Lnf$ of
the committed report.}{Committed report: $\lambda^{*} = 20.416$, mean half-width 0.3855, calibrated
width $2 \times 20.416 \times 0.3855 = 15.74$. Class: $\lambda^{*} = 1.000$, mean half-width 1.5361,
calibrated width 3.072. $\Lnf = 15.74 / 3.072 = 5.123$, over 7,180 replicates of 359 queries.}

\begin{antes}
\antesq[testemunha]{Of the nine certifiable declared queries, how many break under some retraction
of up to three claims, and how many of the seventeen fitted-coefficient queries break at the first?}
\antesq[r0-binario]{Why did the pre-registered Figure~1 rule return no candidate?}
\antesq[nulidade]{Can a query have $\rval = 1$ and $\rnul$ ``never''?}
\end{antes}

\section{T3 in full: a cliff, not a curve}

T3 has 26 rows under the recovering set, estimates at $n = 20\,000$ and seed 0: the nine declared
queries with a valid status and 17 frame queries on the four fitted-coefficient networks. Three declared
queries are excluded with their reasons (\texttt{M-bias}: the published diagram is not a DAG;
\texttt{Schipf\_2010}, TT on T2DM: the outcome is not a possible descendant; \texttt{confounding}: the
treatment has no undirected edge). Read twice, the table says four things.

\textbf{The published analyses are robust and the sampled ones are not.} Eight of the nine declared
queries are unbroken by any retraction to depth three; the exception is \texttt{mediator}, a four-node
textbook diagram. On the fitted networks, 16 of 17 queries have $\rval = 1$. The mechanism is the frame:
a pair is admitted beside an undirected edge, where knowledge can matter, so the measured population
is close to the complement of the queries the source papers were written to answer. That is an
applicability statement for the paper's evaluation section, not a result.

\textbf{The sensitivity curve is a single step.} On 16 of T3's 17 informative rows (the 16 informative
fitted queries and \texttt{mediator}) the interval after one retraction equals the whole-class interval
to four decimals. The exception is again
\texttt{mediator}, whose first retraction covers 57\,\% of the class and whose second saturates it. All
the identifying power the knowledge supplies sits in one claim, and withdrawing it returns the analyst
to where they would be with no knowledge. This is why the pre-registered Figure~1 rule (both radii in
$\{2, 3\}$ and a plateau before the break) found nothing on real coefficients: the shape it looked for
is not in the data, and the declared and applied rows that did have radii 2 or 3 widened 27 to 34 times
before breaking. The companion's recommended figure (its candidate F-A) draws this cliff; the
version on the full eligible pool is in the next section.

\textbf{The null radius and the validity radius are different numbers.} On \texttt{mediator} the set
fails after one retraction and the interval reaches zero after two. On \texttt{ecoli70} eutG $\to$
lacY and both \texttt{magic-niab} G1217 queries, $\rval = 1$ and $\rnul$ is ``never'': the set breaks,
the sign does not. On several fitted rows the band at seed 0 contains zero at budget zero (0 from the
band, 1 from the population): no conclusion to break at that sample size. A sentence about $\rnul$
must name its reading (population, one seed, range over twenty), and a figure's axis label must say
whether it shows conclusions or validity.

\textbf{Why the one- and two-revision columns look identical.} $\Ir{r}$ grows with $r$ and is capped by
the class interval. Of the 303 committed queries under the recovering set, 281 (92.7\,\%) print the
same interval at both budgets: on 76 the class interval is a point, on 205 the first revision already
returns the class. The second column carries information on 20 queries, on \texttt{insurance},
\texttt{Polzer\_2012}, \texttt{Kampen\_2014}, \texttt{sachs}, \texttt{Schipf\_2010} and
\texttt{mediator}, all but the last excluded from T3 by its real-coefficient rule. Replace the
two-revision column with the class interval, and the identity becomes visible as the finding.

\prompt[conceito=r0-binario,dificuldade=0.75]{State the one-claim cliff as measured on T3, name the
exception, and say what it did to the pre-registered Figure~1 rule.}{On 16 of the 17 informative T3
rows the interval after one retraction equals the whole-class interval to four decimals; the exception
is \texttt{mediator}, where one retraction covers 57\,\% of the class and two saturate it. The
pre-registered rule wanted both radii in $\{2, 3\}$ and a plateau before the break on a real-coefficient
network under truthful knowledge; the curve is a single step, so no row qualified, and the rows with
radii 2 or 3 widened 27 to 34 times before breaking.}

\begin{antes}
\antesq[pre-registro]{T3 prints 17 of the 294 eligible sampled queries. By which two rules, and which
one is indefensible?}
\antesq[pre-registro]{What share of printed, dropped and eligible queries has a finite null radius?}
\antesq[testemunha]{Which claim of T3 survives the selection audit?}
\end{antes}

\section{A selection bias in our own table}

T3 is a selection made by two rules. The funnel is the whole path under the recovering set; the
counts are in \texttt{results/ledger/TABLE0\_FUNNEL.md} or were recounted from
\texttt{results/interval/rows.csv} and \texttt{profiles.jsonl} by
\nolinkurl{rho_breakdown/panelB_T3_2026-09-16/make_exhaustive.py}, which first checks the 26 printed
rows cell by cell against \texttt{table2.json}.

\begin{center}\footnotesize
\begin{tabular}{@{}rp{6.2cm}p{6.6cm}@{}}
\toprule
count & stage & what leaves, and why \\
\midrule
39 / 33 / 27 & networks parsed / with a frame / swept & over 500 nodes, or every component over the questionnaire cap \\
80{,}132 & candidate pairs, from observables only & \\
648 & pairs sampled (at most 20 per network) & \\
537 & rows: 528 frame rows in the swept networks + 9 declared & \\
303 & commit to an adjustment set (9 declared, 294 sampled) & 234: on the analyst's graph the outcome cannot descend from the treatment \\
37 & \textbf{rule 1}: fitted coefficients only & 257 on the 23 semi-synthetic networks; \textcolor{verdesc}{defensible}, not in the caption \\
17 & \textbf{rule 2}: finite $\rnul$ first, then frame order, five per network & 20; sorts on the outcome it illustrates: \textcolor{vermesc}{indefensible} \\
\midrule
26 & T3 = 9 declared + 17 sampled & \\
\bottomrule
\end{tabular}
\end{center}

\textbf{Rule 1 is defensible}: on the 23 semi-synthetic networks the effect size is one we drew, and a
named report would present it as a real analysis. It belongs in the caption.

\textbf{Rule 2 is not.} The Table~2 block of \texttt{experiments/interval\_tables.py} (line 1136) sorts the
eligible queries of each network by \texttt{(r0\_pop == CENSORED, frame\_index)} and keeps the first five.
A finite null radius is exactly what the table exists to illustrate.

\begin{medido}[The bias, \texttt{rho\_breakdown/panelB\_T3\_2026-09-16/make\_figs.py} (rerun 16 September)]
Queries on the four fitted networks with a committed set: 37 eligible, 17 printed, 20 dropped (9 on
\texttt{magic-niab}, 8 on \texttt{magic-irri}, 3 on \texttt{ecoli70}). Finite population null radius:
13 of 17 printed (76\,\%), 3 of 20 dropped (15\,\%), 16 of 37 eligible (43\,\%).
\end{medido}

I have to be exact about whose defect this is. The ordering was not chosen after seeing the answer: it
is written in \S6 of \texttt{results/interval/PREREGISTRATION.md} (``finite r0 first, then frame
order''), before any row existed, and the header of \texttt{TABLE2\_REPORTS.md} states it. It was still
the wrong rule, because it sorts on the outcome, and nobody measured its effect until this morning. The
$\rnul$ column of T3 is not representative of the population it was drawn from.

\begin{noprato}[The fix, one line]
In \texttt{experiments/interval\_tables.py}, sort the candidates by \texttt{frame\_index} alone (drop
the \texttt{r0\_pop == CENSORED} term), or better, print all 37 eligible queries so the inclusion rule
has nothing in it a reader must trust. Put rule 1 and the eligible count in the caption. Draw the
figure on the 37, not on the 17.
\end{noprato}

\textbf{What the selection did not bias.} The cliff survives. On all 26 informative queries of the 37
eligible, one retraction returns the whole-class interval (16 of 16 printed, 10 of 10 dropped). Over
all 294 committed sampled queries on 27 networks, 226 are informative; on 205 the one-retraction
interval equals the class interval to four decimals, and on 211 (93\,\%) it covers more than 90\,\% of
the class width. \figref{fig:t3-pool} draws F-A on the 37 eligible queries, printed ones marked,
including the 11 rows where nothing moves, which T3 hides.

\begin{figure}[!ht]
\centering
\includegraphics[width=0.60\linewidth]{t3_f11_pool_dumbbell.pdf}
\caption{All 37 eligible fitted-coefficient queries, printed ones marked with a square: the cliff holds
on the dropped queries too, and the rows with no bar (the class already pins the effect) are the honest
counterweight the printed table leaves out.}
\label{fig:t3-pool}
\end{figure}

\prompt[conceito=pre-registro,dificuldade=0.75]{Give the three finite-null-radius shares of the T3
audit with their denominators, the line that produced the bias, and the fix.}{Finite population
$\rnul$ on 13 of 17 printed queries (76\,\%), 3 of 20 dropped (15\,\%), 16 of 37 eligible (43\,\%). The
sort key \texttt{(r0\_pop == CENSORED, frame\_index)} in \nolinkurl{experiments/interval_tables.py} puts
finite null radii first before taking five per network; the rule is in \S6 of the pre-registration.
Fix: sort by \texttt{frame\_index} alone, or print all 37, and state the real-coefficient rule and the
eligible count in the caption.}

\begin{antes}
\antesq[dois-paineis]{Which candidate does the dossier recommend as Table~1, and what are its two
panels?}
\antesq[decisoes-abertas]{What blocks the abstract of 18 September, according to the dossier?}
\end{antes}

\section{The Table 1 dossier}

The dossier states eleven live claims before any table, each with evidence, significance, nearest
prior art and falsifier (\nolinkurl{rho_breakdown/table1_dossier_2026-09-16/sections/02_claims.tex}).

\begin{center}\footnotesize
\begin{tabular}{@{}lll@{}}
\toprule
& claim & what carries it \\
\midrule
C1 & nobody prices a structural assumption & framing; \textbf{novelty search never run in this wording} \\
C2 & computable from the analyst's inputs & monotonicity audit: 5,000 pairs, 0 violations (bound 0.06\,\%) \\
C3 & claims are the safe unit, hops are not & 3,014 certifiable rows: 0 over-certifications against 39 \\
C4 & the price of a wrong claim & $\Ir{1}$ at 0.851 of the class, premium at zero (Ch.~\ref{ch:gamma}) \\
C5 & the tie with the free screen & $+0.000$ at matched depth (Ch.~\ref{ch:m1b}) \\
C6 & refusal where knowledge matters & 0.910 of 134 committed informative queries \\
C7 & locality in claims, not hops & primary supplier: hop reading above claim depth on 21 of 157 \\
C8 & applicability, as a limit & 8 of 9 declared queries unbroken to depth 3 \\
C9 & the class must be correct & estimated CPDAG: over-certification on 14 of 38 (Ch.~\ref{ch:pest}) \\
C10 & suppliers are an error process & invalid set 0.298 at 72B, 0.288 at 7B \\
C11 & ranking by survival & your $\tau_b$ table; screen column owed \\
\bottomrule
\end{tabular}
\end{center}

\textbf{Thirteen candidates} (T1 to T13, with a variant T1b) are typeset with the numbers that exist and
scored against three requirements (readable after the abstract; verifies the claims we make; compares
with the literature inside the float) and six constraints each earlier design paid
for: one calibrated metric, one population per panel, a row that loses where it must, no discrimination
column, no model family as subject, nothing internal in the float. \textbf{The recommendation is T11},
one float with two panels on two named populations. Its first panel is the certificate ladder over the
3,014 certifiable rows of eleven knowledge conditions on 27 networks (902 committed sets invalid at the
truth): counting claims asserted over-certifies on 729, the separation rule on 15, the hop radius on 39,
the claim radius on 0. Its second panel is the price table at $n = 1000$ with the leave-one-out screen as
a row, and the refusal rate (0.910 of 134) under the float with its own denominator. T3 becomes Table~2,
the policy table T8 becomes Table~3 or prose, your ranking table T4 stays in its section with the screen
column added, and the decision table T5 is the best candidate for the first figure.

Two defects were found and fixed while the dossier was written, each the kind it warns against. The
first draft of T8 and of T11's second panel filled the two-wrong-claim cells of the screen and
stop-rule rows from M1b's 59 committed queries beside cells from the price table's 131: two populations
in one row. And the ladder's 15 over-certifications had been attributed to the leave-one-out screen;
they belong to the separation rule (the screen never over-certifies relative to the stop rule, 0
exceptions in 2,212 rows). A reconciliation, not a defect: my review notes quoted the depth-one screen
gap as $+0.305$, the dossier as $+0.256$; both are in \texttt{TABLE1\_L95.md}, the population line and
the $n = 1000$ line, and every quotation now carries its sample size.

\begin{noprato}[Blocking before the abstract on 18 September]
\textbf{C1's novelty search has never been run in the wording the claim uses.} The searches on record
were pointed at latent-causal representation learning and at Taeb, Guo and Henckel, not at the
sensitivity-analysis lineage C1 says does not already price this, and the statement that Guo and Perkovi\'c never
perturb the knowledge rests on a paper I have not read in full. Either the search runs before the
abstract, or C1 is narrowed to what was searched and the abstract says the narrower thing.
\textbf{The author list must be settled by 18 September}; Chapter~\ref{ch:merge} quotes the rule.
\end{noprato}

\prompt[conceito=dois-paineis,dificuldade=0.7]{Which Table~1 does the dossier recommend, what are its
two panels with their populations, and what blocks the abstract of 18 September?}{T11. Panel one is the
certificate ladder over 3,014 certifiable rows (eleven knowledge conditions, 27 networks, 902 sets invalid
at the truth): the claim count over-certifies on 729, the separation rule on 15, the hop radius on 39,
the claim radius on 0. Panel two is the price table on M1's Panel~A at $n = 1000$ (391, 359, 131
informative queries) with the leave-one-out screen as a row, and the refusal rate 0.910 of 134 printed
under the float with its own denominator. Blocking: the novelty search for C1 has never been run in the
wording of the claim, and the author list must be settled by 18 September.}

\naopasse[pre-registro]{Name the sort key that biased T3, say what it sorts on, and explain in one
sentence why it biases the $\rnul$ column but not the claim that one retraction returns the class
interval. If you cannot, reread the paragraph ``What the selection did not bias''.}

\subsection*{Exercises}

\begin{exercicio}[metodo=table-reading,conceito=nulidade]
Five rows of T3 (recovering set; bands at $n = 20\,000$, seed 0; class interval from the population;
$\rnul$ as population / band at seed 0 / range over 20 seeds).

\smallskip
{\footnotesize\setlength{\tabcolsep}{3pt}
\begin{tabular}{@{}lrllllc@{}}
\toprule
query & true & one revision & two revisions & class & $\rnul$ & $\rval$ \\
\midrule
mediator, X $\to$ Y & $-2.603$ & $[-2.623, -1.109]$ & $[-2.623, 0.000]$ & $[-2.603, 0.000]$ & 2 / 2 / 2 & 1 \\
ecoli70, eutG $\to$ lacY & $0.309$ & $[0.277, 0.377]$ & $[0.277, 0.377]$ & $[0.309, 0.354]$ & never / $>$3 / $>$3 & 1 \\
magic-niab, G418 $\to$ FUS & $0.003$ & $[-0.022, 0.011]$ & $[-0.022, 0.011]$ & $[0.000, 0.003]$ & 1 / 0 / 0..1 & 1 \\
ecoli70, pspB $\to$ nmpC & $-0.108$ & $[-0.531, 0.026]$ & $[-0.531, 0.026]$ & $[-0.525, 0.000]$ & 1 / 1 / 1..3 & 1 \\
arth150, 81 $\to$ 464 & $0.256$ & $[0.097, 0.265]$ & $[0.097, 0.265]$ & $[0.256, 0.256]$ & $>$3 / $>$3 / $>$3 & $\ge 2$ \\
\bottomrule
\end{tabular}}

\smallskip
(a) On which row do the two revision columns differ, and what does that say about the cliff? (b) Which
row has $\rval = 1$ and $\rnul$ ``never'', and why is that not a contradiction? (c) On which row does the
null radius depend on the reading, and how should a sentence quote it? (d) Which row must be left out of
``one retraction returns the class interval'', and why?
\end{exercicio}
\begin{solucao}
(a) \texttt{mediator}: one revision reaches $-1.109$, two reach zero, the class. It is the only row
where the first retraction stops short of the class (57\,\%), so the only one with a second step. (b)
eutG $\to$ lacY. $\rval$ is about the set (one retraction makes $Z$ invalid); $\rnul$ is about the sign
(every graph of the class gives an effect in $[0.309, 0.354]$). Validity breaks, the conclusion does
not. (c) G418 $\to$ FUS: population 1, seed 0 gives 0, range 0 to 1. The true effect is 0.003, so at
$n = 20\,000$ the band usually contains zero before any retraction; quote ``population null radius 1''
or the range, never a bare ``1''. pspB $\to$ nmpC (1 to 3 over seeds) needs the same care. (d) arth150,
81 $\to$ 464: the class interval is a point, so the query is not informative; its band is wide from
sampling alone.
\end{solucao}

\begin{exercicio}[metodo=rewriting,conceito=pre-registro]
T3's current caption reads, in part: \emph{``Sensitivity reports on twenty-six named analyses. \dots
$r_0$ is the null radius, the number of revisions after which the interval contains zero \dots
$r_{\mathrm{val}}$ is the breakdown radius \dots Knowledge: the recovering set; estimates at
$n = 20\,000$, first seed.''} Rewrite it as it should be printed once the fix is applied (all 37
eligible sampled queries shown), so that a reviewer can check the inclusion rule, the counts and the
two radii without opening any file. No internal vocabulary.
\end{exercicio}
\begin{solucao}
\emph{``Sensitivity reports on named analyses. Above the line, the nine published analyses whose
authors declared a treatment and an outcome (three more are excluded: one diagram is not acyclic, in one
the outcome cannot descend from the treatment, in one no claim bears on the treatment); structure as
published, parameters drawn. Below, all 37 sampled queries on the four networks with coefficients fitted
to real data whose analyst graph commits to an adjustment set; networks with drawn parameters are not
shown, since their effect sizes would be ours. Intervals: every effect identified after retracting at
most one claim, and over the whole equivalence class. Null radius: retracted claims until the 95\,\%
band contains zero (population, and range over twenty samples). Validity radius: retracted claims until
the adjustment set is no longer valid. Estimates at $n = 20\,000$.''} What changed: both inclusion rules
and the eligible count are stated, nothing is selected on the outcome, the class interval replaces the
two-revision column, the null radius carries its reading, and ``breakdown radius'' is renamed so it
cannot be confused with the null radius.
\end{solucao}

\section*{Chapter map}
\begin{center}
\begin{tikzpicture}[node distance=0.6cm and 1.1cm]
\node[cmapc] (cell) {a cell: $\lambda^{*}$, calibrated\\ width, ratio to class};
\node[cmap, right=of cell] (t3) {T3: 26 named\\ reports};
\node[cmapg, right=of t3] (cliff) {cliff: one retraction\\ returns the class};
\node[cmapr, below=of t3] (bias) {rule 2 sorts on $\rnul$:\\ 76\,\% vs 43\,\% vs 15\,\%};
\node[cmapg, below=of cliff] (pool) {not biased: 26/26,\\ 211 of 226};
\node[cmap, below=of cell] (fix) {fix: frame order\\ or all 37};
\node[cmapc, right=of pool] (t11) {dossier: T11;\\ C1 search blocks};
\draw[cedge] (cell) -- node[above]{prices} (t3);
\draw[cedge] (t3) -- node[above]{shows} (cliff);
\draw[cedge] (t3) -- node[left]{selected by} (bias);
\draw[cedge] (bias) -- node[above]{repaired by} (fix);
\draw[cedge] (cliff) -- node[right]{audited on} (pool);
\draw[cedge] (pool) -- node[above]{feeds} (t11);
\end{tikzpicture}
\end{center}
